% ==============================================================================
% Appendix
% ==============================================================================
\appendix

\section{Robustness Checks}
\label{app:robustness}

\subsection{Temporal Stability of Local Moran's I}

To assess the temporal stability of the spatial structure, we compute Spearman rank correlations of municipality-level local Moran's $I_i$ values between the June 2024 and June 2025 cross-sections for each status$\times$sex combination.
All correlations are statistically significant ($p < 0.001$).
Values range from 0.227 (Located Dead, male) to 0.607 (Total, total), indicating moderate to strong rank-order stability across the twelve-month interval.
Located Dead shows the weakest correlations (0.227--0.303), consistent with the greater volatility expected from sparse count data.

\begin{table}[htbp]
  \centering
  \caption{Spearman rank correlation of local Moran's $I_i$ between June 2024 and June 2025.}
  \label{tab:spearman}
  \begin{tabular}{llcc}
    \toprule
    Status & Sex & Spearman $\rho$ & $p$-value \\
    \midrule
    Total & total & 0.607 & $<$0.001 \\
    Total & male & 0.526 & $<$0.001 \\
    Total & female & 0.485 & $<$0.001 \\
    Not Located & total & 0.334 & $<$0.001 \\
    Not Located & male & 0.305 & $<$0.001 \\
    Not Located & female & 0.355 & $<$0.001 \\
    Located Alive & total & 0.540 & $<$0.001 \\
    Located Alive & male & 0.494 & $<$0.001 \\
    Located Alive & female & 0.434 & $<$0.001 \\
    Located Dead & total & 0.249 & $<$0.001 \\
    Located Dead & male & 0.227 & $<$0.001 \\
    Located Dead & female & 0.303 & $<$0.001 \\
    \bottomrule
  \end{tabular}
\end{table}

\subsection{Significance Threshold Sensitivity}

Table~\ref{tab:threshold} reports HH municipality counts at three significance thresholds to assess sensitivity to the $\alpha = 0.05$ choice.
At $\alpha = 0.01$, 33--49\% of HH municipalities are retained, confirming that the core clusters remain robust under more stringent thresholds.
No HH municipalities survive $\alpha = 0.001$ with 999 permutations, as expected given the discrete permutation distribution.

\begin{table}[htbp]
  \centering
  \caption{HH municipality counts at different significance thresholds, \referencemonth{} (total sex).}
  \label{tab:threshold}
  \begin{tabular}{lcccc}
    \toprule
    Status & $\alpha=0.05$ & $\alpha=0.01$ & $\alpha=0.001$ & Retention (0.01/0.05) \\
    \midrule
    Total & 106 & 48 & 0 & 45\% \\
    Not Located & 93 & 46 & 0 & 49\% \\
    Located Alive & 110 & 36 & 0 & 33\% \\
    Located Dead & 35 & 12 & 0 & 34\% \\
    \bottomrule
  \end{tabular}
\end{table}

\subsection{Spatial Weights Specification}

All analyses use Queen contiguity weights constructed via \texttt{fuzzy\_contiguity} with a 50-metre buffer on EPSG:6372 projected coordinates.
This approach yields \nislands{} island municipalities.
A prior implementation using exact topological intersection (\texttt{Queen.from\_dataframe}) produced 52 false islands due to sub-metre geometry precision gaps in the source geoparquet---a known issue with high-resolution municipal boundary files.
The 50-metre buffer resolves all boundary gaps without creating spurious connections between non-adjacent municipalities.

\subsection{Jaccard Stability Across Reference Periods}

The Jaccard matrix is broadly stable across reference periods.
In June 2015, NL$\leftrightarrow$LD Jaccard was 0.103 and LA$\leftrightarrow$LD was 0.000; in June 2020, these were 0.164 and 0.198; and in June 2025, 0.143 and 0.151, respectively.
The consistently low values for Located Dead pairs across all three cross-sections reinforce the non-interchangeability finding.
